%
% 01_Einleitung.tex -- Historischer Kontext, die drei Gesetze und
% Zielsetzung der Arbeit
%
% !TEX root = ../../buch.tex
% !TEX encoding = UTF-8
%
\section{Die keplerschen Gesetze
\label{kepler:section:einleitung}}
\kopfrechts{Die keplerschen Gesetze}%

Johannes Kepler formulierte zu Beginn des 17.~Jahrhunderts drei Gesetze über die Bewegung der Planeten um die Sonne.
\index{Kepler, Johannes}%
\index{Sonne}%
\index{Planet}%
Er stützte sich dabei auf die Beobachtungsdaten des dänischen Astronomen Tycho Brahe.
Diese Daten waren für die damalige Zeit aussergewöhnlich präzis.
\index{Brahe, Tycho}%
Die ersten beiden Gesetze veröffentlichte Kepler 1609 in der Astronomia Nova \cite{kepler:astronomianova}.
\index{Astronomia Nova}%
Das dritte Gesetz fand er erst zehn Jahre später.
Er publizierte es 1619 in Harmonices Mundi \cite{kepler:harmonices}.
\index{Harmonices Mundi}%
Die drei Gesetze lauten:
\begin{enumerate}
\item
Die Planeten bewegen sich auf elliptischen Bahnen.
In einem gemeinsamen Brennpunkt dieser Ellipsen steht die Sonne.
\index{Ellipse}%
\item
Ein von der Sonne zum Planeten gezogener Fahrstrahl überstreicht in gleichen Zeiten gleich grosse Flächen.
\index{Fahrstrahl}%
\index{Flachensatz@Flächensatz}%
\item
Die Quadrate der Umlaufzeiten zweier Planeten verhalten sich zueinander wie die Kuben der grossen Halbachsen ihrer Bahnellipsen.
\end{enumerate}
Kepler gewann diese Gesetze rein empirisch aus den Beobachtungsdaten.
Eine physikalische Begründung gelang erst Isaac Newton.
\index{Newton, Isaac}%
Er zeigte 1687 in seinen Principia \cite{kepler:principia}, dass alle drei Gesetze aus dem Gravitationsgesetz folgen.
\index{Principa}%
\index{Graviationsgesetz}%

Diese Arbeit leitet das dritte Gesetz mathematisch aus dem Gravitationsgesetz her.
Die ersten beiden Gesetze beschreiben die Bahn eines einzelnen Planeten.
Das dritte Gesetz vergleicht dagegen verschiedene Planeten miteinander.
Es besagt: Der Quotient $T^2/a^3$ hat für alle Körper um dasselbe Zentralgestirn denselben Wert.
Dabei ist $T$ die Umlaufzeit und $a$ die grosse Halbachse der Bahn.
\index{Umlaufzeit}%
In Lehrbüchern wird das Gesetz meist nur für Kreisbahnen hergeleitet.
Für Kreisbahnen genügt nämlich ein einfaches Kräftegleichgewicht.
Reale Planetenbahnen sind aber Ellipsen.
Für Ellipsen führt die Berechnung der Umlaufzeit auf ein reelles Integral.
Dieses Integral lässt sich mit elementaren Mitteln nur mühsam auswerten.
Elementare Zugänge ohne komplexe Analysis finden sich zum Beispiel bei Vogt \cite{kepler:vogt}.

Wir gehen einen anderen Weg und lösen das Integral mit komplexer Integration.
Die nötigen Methoden wurden im ersten Teil dieses Buches entwickelt.
Abschnitt~\ref{kepler:section:grundlagen} stellt zuerst die geometrischen und physikalischen Grundlagen bereit.
Dort leiten wir die Umlaufzeit als Integral über den Polarwinkel her.
Abschnitt~\ref{kepler:section:kreisbahn} behandelt den Spezialfall der Kreisbahn.
Er dient später als Kontrolle für das allgemeine Resultat.
In Abschnitt~\ref{kepler:section:komplex} verwandeln wir das Integral in ein Wegintegral über den komplexen Einheitskreis.
Dieses Wegintegral werten wir mit der Cauchy-Formel für Ableitungen aus.
Zum Schluss prüfen wir das Resultat an realen Bahndaten von Planeten und einem Kometen.
